\documentclass[11pt,letterpaper]{article}
\usepackage[margin=1in]{geometry}
\usepackage{setspace}
\usepackage{enumitem}
\usepackage{array}
\usepackage{booktabs}
\usepackage{tabularx}
\usepackage{graphicx}
\usepackage{titlesec}
\usepackage{fancyhdr}
\usepackage[T1]{fontenc}
\usepackage[utf8]{inputenc}
\usepackage{lmodern}

\setstretch{1.03}
\pagestyle{empty}
\setlength{\parindent}{15pt}
\setlength{\parskip}{4pt}
\setlist{nosep,leftmargin=0.24in}
\titlespacing*{\section}{0pt}{0.35em}{0.1em}
\titlespacing*{\subsection}{0pt}{0.25em}{0.08em}
\titleformat{\section}{\normalfont\bfseries\large}{}{0pt}{}
\titleformat{\subsection}{\normalfont\bfseries}{}{0pt}{}

\begin{document}
\begin{center}
{\Large\bfseries Budget Justification}\\[0.05in]
{\normalsize Mostafa Rezaali}\\
{\normalsize NSF AGS-PRF Application}
\end{center}
\vspace{-0.08in}
\section*{Overview}
This AGS-PRF fellowship budget follows the solicitation-specified fellowship structure. The budget contains stipend support for the Fellow and a fellowship allowance. No voluntary committed cost sharing is included. The requested funds are mapped to the standard AGS-PRF structure and are justified by the research, dissemination, data-management, and professional-development activities described in the Project Description.

\section*{Year 1 Stipend}
\textbf{Stipend: \$70,000.} The stipend will support full-time fellowship effort by Mostafa Rezaali on research, broader impacts, mentoring, dissemination, and professional development. Activities include developing the conditional-flow forecast model, preparing the cross-validated hindcast archive, computing deterministic and probabilistic verification metrics, documenting workflows, and preparing the first manuscript. The stipend is not requested as salary for institutional personnel; it supports the Fellow's full-time training and research effort.

\section*{Year 1 Fellowship Allowance}
\textbf{Fellowship allowance: \$30,000.} Planned uses include health insurance and benefit-related costs; high-performance computing, cloud storage, and secure data-storage expenses; workstation or peripheral support needed for model development and analysis; conference and workshop travel to present results and receive feedback; open-access publication and data-archive fees; and modest software, training, and dissemination costs directly related to the fellowship. These costs are necessary because probabilistic graph-based climate forecasting requires repeated model training, ensemble evaluation, large derived data products, and transparent public release of results.

\section*{Year 2 Stipend}
\textbf{Stipend: \$72,000.} The second-year stipend will support full-time fellowship effort including probabilistic ensemble training, uncertainty-calibration analyses, physical-regime diagnostics, flash-drought-relevant risk evaluation, manuscript preparation, open-science release, student mentoring, and career-development activities. The second year emphasizes completion, interpretation, dissemination, and transition from model development to independent research leadership.

\section*{Year 2 Fellowship Allowance}
\textbf{Fellowship allowance: \$30,000.} Planned uses mirror Year 1 and will support continued health/benefit costs, computing and storage, conference participation, publication and data dissemination, project-specific supplies, software, and professional development. Anticipated travel includes at least one national conference such as AGU, AMS, or a comparable climate/AI meeting, plus targeted workshops or short training opportunities in subseasonal prediction, uncertainty quantification, or responsible AI.

\section*{Computing, Storage, and Software Rationale}
The project requires GPU-enabled training, repeated cross-validation, ensemble generation, reliability analysis, and storage of gridded hindcast products. Allowance funds may support computing allocations, cloud or institutional storage, backup media, data-transfer charges, and workstation peripherals needed to maintain efficient and reproducible workflows. Software costs, if any, will be limited to project-specific analysis, visualization, archive preparation, or accessibility needs not already provided by the host institution.

\section*{Travel, Publication, and Dissemination Rationale}
Travel funds will allow the Fellow to present results, receive critique from the atmospheric-science community, build professional networks, and identify pathways for translating probabilistic heat-risk diagnostics into useful climate services. Publication and archive funds will support open-access manuscripts, persistent data/code releases, and long-term preservation of reproducible research products.

\section*{Budget Control and Compliance}
All spending will be tied directly to fellowship research, training, or dissemination objectives. The Fellow will track allowance spending, maintain documentation for project-related expenses, and follow NSF, Research.gov, and host-institution rules. No funds are requested for unrelated personnel, construction, entertainment, or voluntary committed cost sharing.

\newpage
\section*{Detailed Computing and Storage Justification}
The model-development plan requires repeated training and evaluation of deterministic and probabilistic graph models across leave-years-out validation folds. The fellowship allowance may support GPU-enabled computing, scratch storage, archival storage, data-transfer costs, and backup resources needed to preserve reproducible products. These costs are tied directly to the project because CFM ensembles require multiple model samples, verification files, and derived diagnostics for reliability, CRPS, and regional skill analysis.

Storage needs will include processed PRISM and ERA5-derived predictors, fold-specific model inputs, trained model checkpoints, generated ensemble hindcasts, skill maps, summary tables, and manuscript figures. Funds may also support secure storage and backup media to reduce the risk of data loss during the fellowship.

\newpage
\section*{Detailed Travel and Dissemination Justification}
Travel funds will support presentation of project results at atmospheric-science, climate, geospatial, or AI-for-science meetings. These venues are important for receiving technical feedback, comparing model diagnostics with related subseasonal prediction work, building collaborations, and developing the Fellow's independent professional network. Allowable costs may include registration, transportation, lodging, meals, and related travel expenses consistent with institutional and NSF rules.

Dissemination costs may include open-access publication fees, data-archive costs, persistent DOI creation, accessibility preparation, and repository documentation. These costs are justified because the project explicitly commits to transparent release of nonrestricted derived products, code, and teaching materials.

\newpage
\section*{Detailed Training and Professional Development Justification}
The fellowship allowance may support short courses, workshops, or professional-development activities in subseasonal prediction, uncertainty quantification, scientific software, responsible AI, or climate-risk communication. These activities will strengthen the Fellow's ability to lead an independent research program and to translate AI methods into atmospheric-science questions.

Professional-development spending will be modest and project-relevant. Priority will be given to activities that improve research quality, reproducibility, mentoring capacity, or dissemination. The requested budget therefore supports both the scientific execution of the project and the AGS-PRF goal of developing an independent early-career scientist.
\end{document}
